\chapter{Prompt Engineering Fundamentals\\หลักการใช้ AI เพื่อยกระดับการ Mentoring}

AI สามารถช่วยจัดโครงสร้างข้อมูล เปรียบเทียบทางเลือก และสร้างร่างแรกได้รวดเร็ว แต่ไม่รู้จัก Mentee เท่ากับบทสนทนาจริง ไม่รับผิดชอบต่อผลของคำแนะนำ และอาจสร้างข้อมูลที่ดูน่าเชื่อถือแต่ไม่จริง บทบาทที่เหมาะสมจึงเป็น \textbf{ผู้เร่งการคิด} มิใช่ผู้แทนการตัดสินใจของ Mentor หรือ Mentee

\learningobjectives{
  \item เขียน Prompt ด้วยกรอบ RTCF
  \item ใช้ AI โดยรักษาความลับ ความถูกต้อง และความเป็นเจ้าของ
  \item ออกแบบ workflow ที่มีมนุษย์ตรวจสอบทุกจุดสำคัญ
}

\section{RTCF: ไวยากรณ์ของ Prompt ที่ชัด}
\begin{longtable}{P{.16\textwidth}P{.28\textwidth}P{.42\textwidth}}
\toprule
\textbf{ส่วน} & \textbf{หน้าที่} & \textbf{คำถามออกแบบ}\\
\midrule
Role & กำหนดมุมมองของ AI & ต้องการความเชี่ยวชาญและเกณฑ์แบบใด\\
Task & ระบุงานที่ต้องทำ & ผลลัพธ์ต้องช่วยการตัดสินใจอะไร\\
Context & ให้ข้อมูลและข้อจำกัด & ข้อเท็จจริงใดจำเป็น อะไรห้ามสมมติ\\
Format & กำหนดรูปแบบผลลัพธ์ & ใครจะใช้ และควรอ่านอย่างไร\\
\bottomrule
\end{longtable}

\begin{promptbox}
\textbf{Role:} ท่านเป็นผู้ช่วยวิเคราะห์ Tech Portfolio ที่ระมัดระวังเรื่องหลักฐาน\\
\textbf{Task:} ระบุช่องว่างสำคัญที่สุด 3 ข้อและคำถามที่ Mentor ควรถาม\\
\textbf{Context:} ใช้เฉพาะข้อมูล Portfolio ด้านล่าง ห้ามสร้างข้อมูลเพิ่ม และแยกข้อเท็จจริงออกจากข้อสันนิษฐาน\\
\textbf{Format:} ตาราง 4 คอลัมน์: ช่องว่าง, หลักฐานที่เห็น, สิ่งที่ยังไม่ทราบ, คำถามสำหรับ Mentee
\end{promptbox}

\section{วงจร Human-in-the-Loop}
workflow ที่รับผิดชอบมีห้าจังหวะ: Mentor กำหนดคำถาม; คัดข้อมูลที่แบ่งปันได้; AI สร้างร่างหรือทางเลือก; Mentor และ Mentee ตรวจข้อเท็จจริง ตีความ และแก้ไข; จากนั้นมนุษย์ตัดสินใจและรับผิดชอบผล AI ควรช่วยเพิ่มคุณภาพการสนทนา ไม่ใช่ลดการสนทนา

\section{หลักการใช้ AI อย่างรับผิดชอบ}
\subsection{ความลับและความยินยอม}
อย่านำข้อมูลส่วนบุคคล แนวคิดที่ยังไม่เปิดเผย ข้อมูลวิจัยอ่อนไหว หรือรายละเอียดที่ระบุตัว Mentee เข้าระบบโดยไม่มีสิทธิและความยินยอม ควรลดทอนข้อมูล ใช้นามสมมติ และปฏิบัติตามนโยบายของมหาวิทยาลัยและโครงการ

\subsection{ความถูกต้องและแหล่งอ้างอิง}
กำชับให้ AI ใช้เฉพาะข้อมูลที่ให้ และทำเครื่องหมายสิ่งที่ไม่ทราบ ทุกตัวเลข ข้อเท็จจริง และการอ้างอิงต้องตรวจจากแหล่งต้นทางก่อนใช้ ไม่ควรให้ความลื่นไหลของภาษาแทนความจริง

\subsection{ความเป็นเจ้าของและอคติ}
Mentee ต้องสามารถอธิบายและปกป้องข้อความสุดท้ายได้ AI อาจสะท้อนอคติจากข้อมูลหรือกรอบคำถาม Mentor จึงควรถามว่าใครถูกมองข้าม มุมมองใดไม่ปรากฏ และคำแนะนำนี้ให้ประโยชน์แก่ใคร

\begin{examplebox}
AI เสนอว่า Portfolio ขาด \enquote{ประสบการณ์เชิงพาณิชย์} Mentor ไม่ควรรับข้อสรุปทันที แต่อาจใช้เป็นคำถามว่า \enquote{เป้าหมายของท่านต้องการเส้นทางเชิงพาณิชย์ หรือเส้นทางนโยบาย/สาธารณะมากกว่า} การตรวจเป้าหมายก่อนแก้ Gap ป้องกันไม่ให้ AI กำหนดนิยามความสำเร็จแทน Mentee
\end{examplebox}

\begin{mentornote}
หลังอ่านผลจาก AI ให้ถามสามข้อเสมอ: ข้อใดเป็นข้อเท็จจริง ข้อใดเป็นการตีความ และข้อใดต้องกลับไปถาม Mentee
\end{mentornote}

\begin{keytakeaway}
Prompt ที่ดีช่วยให้ AI ทำงานชัดขึ้น แต่ความรับผิดชอบยังอยู่กับมนุษย์ ใช้ AI เพื่อสร้างทางเลือกและคำถาม ไม่ใช่เพื่อมอบอำนาจการตัดสินใจหรือความสัมพันธ์
\end{keytakeaway}

\begin{reflection}
นำ Prompt ที่ท่านเคยใช้มาปรับด้วย RTCF เพิ่มข้อจำกัดเรื่องข้อมูลและรูปแบบการตรวจสอบ จากนั้นระบุจุดอย่างน้อยสองจุดที่มนุษย์ต้องทบทวนก่อนใช้ผลลัพธ์
\blankline\blankline
\end{reflection}

\chaptersummary{AI มีคุณค่ามากที่สุดเมื่อทำงานภายในกรอบที่ชัด มีข้อมูลที่เหมาะสม และมีมนุษย์ตรวจสอบ RTCF ช่วยออกแบบคำสั่ง ส่วนหลักความลับ ความถูกต้อง ความเป็นเจ้าของ และอคติช่วยรักษาคุณภาพของการ Mentoring}
